%% -*- coding: utf-8 -*-
%% EXAMPLE CONTENT — ships with a new SciTeX workspace. Replace it with yours.
%% Every figure quoted here is regenerated by scripts/reproduce_figures.py.
\section{Results}
\label{sec:results}
\pdfbookmark[1]{Results}{results}

\subsection{The data as the classifier sees it}

Figure \ref{fig:01_digit_grid} shows one example of each digit at the
resolution the model actually receives. At $8\times8$ the strokes are only two
or three pixels wide, and several classes already look alike --- notably 8 and
1, whose ink occupies a similar central column.

\subsection{Held-out accuracy}

The linear support vector machine labels 131 of the 150 held-out images
correctly, an accuracy of 87.3\%. For reference, always guessing a single digit
would score about 10\% on this balanced test set, so nearly all of the
achievable signal is recovered by a linear rule on raw pixels alone.

\subsection{Where the errors are}

Accuracy alone hides the interesting structure. Figure
\ref{fig:02_confusion_matrix} gives the full confusion matrix over the 150
held-out images, 15 per digit. Errors are strongly asymmetric across classes:

\begin{itemize}
  \item The digit 0 is recognised perfectly (15/15) and is never predicted for
    any other class.
  \item Seven further digits --- 2, 3, 4, 5, 6, 7 and, marginally, 1 --- sit at
    13/15 or 14/15, with their misses scattered singly rather than concentrated
    on one rival class.
  \item The digit 9 falls to 12/15, and two of its three errors are predictions
    of 7.
  \item The digit 8 is the clear outlier at 7/15. Its errors are spread over
    five different classes, most often 1 (three cases) and 6 (two cases).
\end{itemize}

Every off-diagonal count in Figure \ref{fig:02_confusion_matrix} is printed in
the cell, so these statements can be read off the figure directly rather than
taken on trust.

%%%% EOF
